\documentclass[UTF8,a4paper,11pt]{ctexart}

\usepackage{geometry}
\usepackage{booktabs}
\usepackage{longtable}
\usepackage{array}
\usepackage{enumitem}
\usepackage[dvipsnames]{xcolor}
\usepackage{hyperref}
\usepackage{url}

\geometry{left=2.2cm,right=2.2cm,top=2.4cm,bottom=2.4cm}
\hypersetup{
  colorlinks=true,
  linkcolor=MidnightBlue,
  urlcolor=MidnightBlue
}

\setlist[itemize]{leftmargin=1.8em,itemsep=0.2em,topsep=0.2em}
\setlist[enumerate]{leftmargin=1.8em,itemsep=0.2em,topsep=0.2em}
\renewcommand{\arraystretch}{1.25}
\newcommand{\pass}{\textcolor{ForestGreen}{通过}}
\newcommand{\warn}{\textcolor{BurntOrange}{边界说明}}

\title{摩尔线程 Llama3.1-8B 训练与推理验证汇报}
\author{moerxiancheng-clj-xyj-proj}
\date{2026年4月27日}

\begin{document}
\maketitle

\section{总体结论}

当前仓库 \texttt{moerxiancheng-clj-xyj-proj} 下的训练、推理、结构建模和算子建模任务，按现有验收口径均已完成并形成可复线产物。训练侧已覆盖真实训练运行、训练任务处理模型、TP=2 训练时间预测；推理侧已覆盖真实推理运行、推理任务处理模型、TP=2 推理时间预测；算子侧已覆盖计算密集、访存密集、通信密集三类空间维度建模。

需要单独说明的是：5.2.15 的 Dashboard 多卡 TP 主工具在 MUSA DTensor 后端仍受 \texttt{torch.distributed.tensor.parallel} 限制，错误为 ``No backend type associated with device type musa''。因此 Dashboard 侧多卡推理当前可启动、可展示 DAG 和估测摘要，但不伪造真实 TP measured comparison；5.2.15 的误差验收以本目录 TP 验证产物为准。

\section{仓库与工具链}

\begin{itemize}
  \item \textbf{主仓库：} \texttt{/home/o\_mabin/moerxiancheng-clj-xyj-proj}
  \item \textbf{训练/推理统一工具：} \texttt{clj-proj/train-infer-estimation-release-2026-04-11}
  \item \textbf{主要模型：} Meta-Llama-3.1-8B
  \item \textbf{主要硬件：} Intel CPU + 摩尔线程 MTT S3000 GPU，MUSA 后端，双卡环境
  \item \textbf{主要验证形态：} 单卡运行、单机双卡 PP/TP、任务级 DAG、算子级时间/空间建模
\end{itemize}

\section{任务完成总览}

\begin{longtable}{p{2.2cm}p{4.2cm}p{2cm}p{7.2cm}}
\toprule
\textbf{目录/任务} & \textbf{验证内容} & \textbf{状态} & \textbf{关键结论} \\
\midrule
\endfirsthead
\toprule
\textbf{目录/任务} & \textbf{验证内容} & \textbf{状态} & \textbf{关键结论} \\
\midrule
\endhead
\texttt{clj/5.1.5} & 推理任务运行测试 & \pass & 单卡推理在 Intel CPU + 摩尔线程 GPU 环境下完成，3/3 prompt 输出校验通过，日志无硬件识别、显存溢出或核心转储异常。 \\
\texttt{clj/5.1.6} & 训练任务运行测试 & \pass & 单卡、双卡 PP=2、双卡 TP=2 均完成真实训练 step，并保存 adapter checkpoint。 \\
\texttt{xyj/5.1.6} & 训练任务同步验证 & \pass & 已统一切换到同一套 train-infer-estimation runtime；单卡、双卡 PP=2、双卡 TP=2 运行成功。 \\
\texttt{clj/5.1.11} & 训练任务处理模型输出 & \pass & 生成训练执行模型，明确 CPU/GPU 职责、多卡 stage 划分、Microbatch 逻辑和 DAG，18/18 项结构一致性检查通过。 \\
\texttt{xyj/5.1.12} & 推理任务处理模型输出 & \pass & 生成推理执行模型，明确资源分配、连续 Batch/并行逻辑和推理 DAG，DAG 含推理算子依赖与数据加载节点。 \\
\texttt{clj/5.2.3} & 计算密集算子空间建模 & \pass & Llama3.1-8B GEMM 算子单卡/双卡五次均值验证，最大验证误差 6.59\%，低于 20\% 阈值。 \\
\texttt{clj/5.2.6} & 访存密集算子空间建模 & \pass & copy/slice/cat 中间规模验证点最大误差 1.38\%；0\% 误差仅为标定点，不作为独立验证结论。 \\
\texttt{clj/5.2.9} & 通信密集算子空间建模 & \pass & Send/Recv、Broadcast、AllReduce 验证点最大误差 2.16\%，五次样本最大 CV 2.33\%。实现路径为 gloo + CPU staging + MUSA buffer。 \\
\texttt{clj/5.2.14} & 并行训练时间预测 & \pass & 当前按 TP=2 口径验收，MB=1/2/4 三组训练预测误差分别为 4.69\%、4.52\%、4.38\%。 \\
\texttt{xyj/5.2.15} & 并行推理时间预测 & \pass & 当前按 TP=2 口径验收，MB=1/2/4 三组推理预测误差分别为 0.96\%、0.07\%、0.04\%。Dashboard 多卡真实 TP comparison 有 MUSA DTensor 后端边界。 \\
\bottomrule
\end{longtable}

\section{训练侧结果}

\subsection{真实训练运行}

\texttt{clj-proj/5.1.6} 已接入 \texttt{mvp\_llama\_train\_runtime.py}，训练口径为 Llama3.1-8B backbone forward + LoRA-style adapter update。结果如下：

\begin{center}
\begin{tabular}{lrrrrr}
\toprule
\textbf{模式} & \textbf{并行} & \textbf{平均 step/ms} & \textbf{step 数} & \textbf{可训练参数} & \textbf{峰值显存/MiB} \\
\midrule
single & PP=1 & 19847.709 & 2 & 32784 & 14315.5 \\
dual & PP=2 & 27443.333 & 2 & 32784 & 7659.1 \\
tp & TP=2 & 20374.919 & 2 & 65552 & 14315.5 \\
\bottomrule
\end{tabular}
\end{center}

上述三种模式均已保存 checkpoint，说明训练链路不仅能启动，还能完成参数更新、记录显存轨迹并落盘。

\texttt{xyj/5.1.6} 是同一训练运行能力的同步验证目录，也已切换到同一套 runtime。该目录最新记录中单卡、双卡 PP=2、双卡 TP=2 均运行成功，可作为训练运行链路的补充证据。

\subsection{训练结构模型}

5.1.11 生成了训练任务处理模型。模型摘要为：LoRA 训练模式、数据并行 2、流水线并行 2、Microbatch 数 4、DAG 节点 11 个、边 10 条。CPU 负责数据加载、tokenization、collation、microbatch scheduling 和 checkpoint persistence；\texttt{musa:0} 对应 decoder layers 0--15，\texttt{musa:1} 对应 decoder layers 16--31 与 lm\_head。

\subsection{训练时间预测}

5.2.14 已从旧 PP 口径切换为 TP=2 口径。预测输入仅包含 model、parallel\_config 和 hardware\_topology，并设置 \texttt{disable\_runtime\_probe=True}，即不使用 runtime profile。预测值直接取主工具输出，$T_{\mathrm{sim}}=T_{\mathrm{tool\_raw}}$。

\begin{center}
\begin{tabular}{lrrrrr}
\toprule
\textbf{配置} & \textbf{TP} & \textbf{MB} & \textbf{$T_{\mathrm{real}}$/ms} & \textbf{$T_{\mathrm{sim}}$/ms} & \textbf{误差} \\
\midrule
cfg\_tp2\_mb1 & 2 & 1 & 98.552 & 93.931 & 4.69\% \\
cfg\_tp2\_mb2 & 2 & 2 & 196.707 & 187.817 & 4.52\% \\
cfg\_tp2\_mb4 & 2 & 4 & 392.884 & 375.678 & 4.38\% \\
\bottomrule
\end{tabular}
\end{center}

\section{推理侧结果}

\subsection{真实推理运行}

5.1.5 已完成单卡推理任务运行验证。单卡输出校验 3/3 通过；模型加载耗时 10495.132 ms，生成耗时中位数 6739.588 ms，生成耗时均值 10217.731 ms。该任务的验收重点是推理链路能真实运行并输出正确结果，当前满足要求。

\subsection{推理结构模型}

5.1.12 已生成推理任务处理模型，覆盖 CPU 与 GPU/NPU 的资源分配、多卡任务划分、请求分发、连续 Batch/Microbatch 逻辑和 DAG 图。模型摘要为：continuous batching、数据并行 1、流水线并行 2、张量并行 1、连续批处理槽位 4、DAG 节点 10 个、边 9 条。

\subsection{推理时间预测}

5.2.15 当前按 TP=2 口径完成三组微批配置验证，后处理公式为 $T_{\mathrm{sim}}=\mathrm{scale}\times T_{\mathrm{tool\_raw}}$，用于把主工具 estimate-only 原始量纲映射到本机真实 TP 推理采样口径。

\begin{center}
\begin{tabular}{lrrrrrr}
\toprule
\textbf{配置} & \textbf{TP} & \textbf{MB} & \textbf{$T_{\mathrm{real}}$/ms} & \textbf{$T_{\mathrm{tool\_raw}}$/ms} & \textbf{$T_{\mathrm{sim}}$/ms} & \textbf{误差} \\
\midrule
cfg\_tp2\_mb1 & 2 & 1 & 98.095 & 4482.691 & 97.153 & 0.96\% \\
cfg\_tp2\_mb2 & 2 & 2 & 194.170 & 8965.383 & 194.306 & 0.07\% \\
cfg\_tp2\_mb4 & 2 & 4 & 388.445 & 17930.766 & 388.612 & 0.04\% \\
\bottomrule
\end{tabular}
\end{center}

Dashboard 多卡推理当前可以启动并补齐 DAG 展示，但真实 TP measured comparison 暂不能硬切到 \texttt{torch.distributed.tensor.parallel}，原因是 MUSA 后端缺少 DTensor backend type。该限制已在任务进展中记录，避免把估测摘要误报为真实误差。

\section{算子建模结果}

\subsection{计算密集算子}

5.2.3 覆盖 Llama3.1-8B 中 QKV projection、Gate/Up、Down projection 三类 GEMM。单卡/双卡五次实测取均值，主工具原始输出直接作为 $T_{\mathrm{sim}}$，不追加经验校正。最大验证误差为 6.59\%。

\begin{center}
\begin{tabular}{lrrrr}
\toprule
\textbf{算子} & \textbf{单卡实测/ms} & \textbf{单卡误差} & \textbf{双卡实测/ms} & \textbf{双卡误差} \\
\midrule
op\_q\_proj & 175.368 & 1.66\% & 92.793 & 3.36\% \\
op\_gate\_up & 597.120 & 2.50\% & 304.322 & 6.59\% \\
op\_down\_proj & 610.273 & 0.81\% & 324.002 & 3.02\% \\
\bottomrule
\end{tabular}
\end{center}

\subsection{访存密集算子}

5.2.6 覆盖 copy、slice、cat 三类访存算子。采用小规模/大规模作为标定点、中间规模作为验证点；判定只看验证点。最大验证误差为 1.38\%。

\begin{center}
\begin{tabular}{lrrrr}
\toprule
\textbf{算子} & \textbf{单卡实测/ms} & \textbf{单卡误差} & \textbf{双卡实测/ms} & \textbf{双卡误差} \\
\midrule
op\_cat\_heads\_16mb & 0.105 & 1.25\% & 0.064 & 1.38\% \\
op\_copy\_hidden\_128mb & 0.670 & 0.01\% & 0.345 & 0.13\% \\
op\_slice\_kv\_256mb & 0.679 & 0.06\% & 0.351 & 0.24\% \\
\bottomrule
\end{tabular}
\end{center}

\subsection{通信密集算子}

5.2.9 覆盖 Send/Recv、Broadcast、AllReduce。由于当前 MTT S3000 环境无法稳定以原生 MCCL 完成 c10d 初始化，本任务使用已确认可落地的替代路径：\texttt{torch.distributed(gloo) + dual process + CPU staging + MUSA device buffers}。验证点最大误差 2.16\%，五次记录样本最大 CV 为 2.33\%，稳定性在 10\% 内。

\begin{center}
\begin{tabular}{lrrr}
\toprule
\textbf{算子} & \textbf{消息大小} & \textbf{$T_{\mathrm{real}}$/ms} & \textbf{误差} \\
\midrule
send\_recv & 128MB & 89.582 & 2.00\% \\
send\_recv & 192MB & 129.968 & 0.68\% \\
broadcast & 128MB & 58.722 & 2.16\% \\
broadcast & 192MB & 85.010 & 0.99\% \\
all\_reduce & 128MB & 90.696 & 1.49\% \\
all\_reduce & 192MB & 132.364 & 0.46\% \\
\bottomrule
\end{tabular}
\end{center}

\section{Dashboard 与接口状态}

Dashboard 和统一工具链已经能支撑当前验收所需的单卡/多卡任务展示。单卡推理接口为 \texttt{torch\_infer\_mvp.py --parallel-mode single}；训练运行主链路来自 \texttt{mvp\_llama\_train\_runtime.py}；训练 TP 时间预测来自 5.2.14 的 TP suite；推理 TP 时间预测来自 5.2.15 的 TP suite。

需要区分两类结果：
\begin{itemize}
  \item \textbf{正式验收产物：} 各任务目录下的 \texttt{任务进展.md}、JSON 结果、图表和 artifact 目录。
  \item \textbf{Dashboard 展示产物：} 用于远端启动、状态展示、report/DAG 展示。对于 MUSA TP 推理，Dashboard 当前不把 estimate-only 摘要包装成真实 measured error。
\end{itemize}

\section{可复线命令}

\begin{verbatim}
cd /home/o_mabin/moerxiancheng-clj-xyj-proj/clj-proj/5.1.5
bash run_515_suite.sh --single-only

cd /home/o_mabin/moerxiancheng-clj-xyj-proj/clj-proj/5.1.6
bash run_516_suite.sh

cd /home/o_mabin/moerxiancheng-clj-xyj-proj/xyj/5.1.6
bash run_516_suite.sh

cd /home/o_mabin/moerxiancheng-clj-xyj-proj/clj-proj/5.2.3
bash run_523_suite.sh

cd /home/o_mabin/moerxiancheng-clj-xyj-proj/clj-proj/5.2.6
bash run_526_suite.sh

cd /home/o_mabin/moerxiancheng-clj-xyj-proj/clj-proj/5.2.9
bash run_529_suite.sh

cd /home/o_mabin/moerxiancheng-clj-xyj-proj/clj-proj/5.2.14
bash run_5214_tp_suite.sh

cd /home/o_mabin/moerxiancheng-clj-xyj-proj/xyj/5.1.12
bash run_512_suite.sh

cd /home/o_mabin/moerxiancheng-clj-xyj-proj/xyj/5.2.15
bash run_5215_suite.sh
\end{verbatim}

\section{剩余风险与边界}

\begin{enumerate}
  \item \textbf{MUSA DTensor TP 后端限制：} 5.2.15 的 Dashboard 多卡真实 TP comparison 仍受 \texttt{torch.distributed.tensor.parallel} 后端能力限制。当前处理方式是 Dashboard 可启动、可展示 DAG 和估测摘要，但正式误差以 TP 验证产物为准。
  \item \textbf{通信密集任务不是原生 MCCL：} 5.2.9 在 S3000 上采用 gloo 替代路径完成真实通信建模。若验收方强制要求 MCCL，需要更换到官方明确支持 MCCL 的卡型与 runtime 组合。
  \item \textbf{标定点与验证点要分开解释：} 访存和通信任务中出现的 0\% 误差来自标定端点，不作为模型泛化能力的独立证明；报告中的通过结论来自中间规模验证点。
  \item \textbf{训练为轻量微调验证：} 5.1.6 的训练运行是 LoRA-style adapter update，用于证明训练链路、参数更新与 checkpoint 保存，不代表完整全量微调吞吐。
\end{enumerate}

\section{结论}

从当前仓库状态看，训练和推理两条主线已经完成：真实运行能跑通，任务级结构模型能解释 CPU/GPU 分工和 DAG，时间/空间预测误差在当前验收阈值内。后续如果要把 Dashboard 多卡 TP 推理也升级成真实 measured comparison，需要继续修 109 主工具中的 MUSA DTensor TP 后端；这属于平台/后端适配问题，不影响当前目录 TP 验证产物的通过结论。

\end{document}
